\documentclass[11pt,a4paper]{article}

% ---------- Idioma y codificación ----------
\usepackage[spanish,es-tabla]{babel}
\usepackage[utf8]{inputenc}
\usepackage[T1]{fontenc}
\usepackage{lmodern}

% ---------- Formato de página ----------
\usepackage[margin=2.5cm]{geometry}
\usepackage{setspace}
\setstretch{1.15}
\usepackage{parskip}

% ---------- Tablas y listas ----------
\usepackage{booktabs}
\usepackage{tabularx}
\usepackage{array}
\usepackage{enumitem}

% ---------- Títulos en negrita (sin color) ----------
\usepackage{titlesec}
\titleformat{\section}
  {\Large\bfseries}
  {\thesection}{0.8em}{}
\titleformat{\subsection}
  {\large\bfseries}
  {\thesubsection}{0.6em}{}
\titleformat{\subsubsection}
  {\normalsize\bfseries}
  {\thesubsubsection}{0.5em}{}
\titlespacing*{\section}{0pt}{18pt}{8pt}

% ---------- Encabezado y pie ----------
\usepackage{fancyhdr}
\pagestyle{fancy}
\fancyhf{}
\fancyhead[L]{\small\textbf{Scrum --- Sprint \arabic{sprintnum}}}
\fancyhead[R]{\small\textbf{CrediStep}}
\fancyfoot[C]{\small\thepage}
\renewcommand{\headrulewidth}{0.4pt}

% ---------- Número de sprint ----------
\newcounter{sprintnum}
\setcounter{sprintnum}{6}   % Cambie este número en cada sprint

% ---------- Caja del objetivo (solo blanco y negro) ----------
\usepackage[most]{tcolorbox}
\newtcolorbox{objetivo}{
  colback=white, colframe=black,
  boxrule=0.8pt, arc=0pt,
  left=8pt, right=8pt, top=6pt, bottom=6pt,
  fonttitle=\bfseries, title={Objetivo del Sprint}
}

% ======================================================
\begin{document}

% ---------- Portada ----------
\begin{center}
  {\Huge\bfseries Documento de Sprint} \\[0.4cm]
  {\LARGE\bfseries Sprint \arabic{sprintnum}} \\[0.3cm]
  {\large CrediStep} \\[1.2cm]
\end{center}

\section{Información general}

\begin{tabularx}{\textwidth}{>{\bfseries}l X}
  \toprule
  Sprint               & \arabic{sprintnum} \\
  Proyecto             & CrediStep \\
  Product Owner        & Víctor Hugo Vela Pinzón (profesor) \\
  Scrum Master         & Beiker Romario Jaramillo Giraldo \\
  Equipo de desarrollo & Laura Sofía Ortiz Gómez, Sarahy Michelle Silva Pinto, Valentina Carrillo Peñuela, Miguel Ángel Gómez López \\
  Fecha de inicio      & 13/10/2026 \\
  Fecha de fin         & 19/10/2026 \\
  Duración             & 1 semana \\
  Capacidad del equipo & 8 puntos \\
  \bottomrule
\end{tabularx}

% ---------- 2. Objetivo del sprint ----------
\section{Objetivo del sprint}

\begin{objetivo}
  Dibujar a mano el flujo principal de CrediStep y presentarlo en el taller práctico de prototipo rápido.
\end{objetivo}

% ---------- 3. Backlog del sprint ----------
\section{Backlog del sprint}

\subsection{Historias de usuario seleccionadas}

\begin{tabularx}{\textwidth}{c X c c c}
  \toprule
  \textbf{ID} & \textbf{Historia de usuario} & \textbf{Prioridad} & \textbf{Puntos} & \textbf{Estado} \\
  \midrule
  HU-01 & Como equipo, quiero dibujar las pantallas de inicio de sesión y diagnóstico, para mostrar cómo empieza el usuario. & Alta & 3 & Terminada \\
  HU-02 & Como equipo, quiero dibujar la pantalla de resultado, el plan y la acción recomendada, para mostrar cómo avanza el usuario. & Alta & 3 & Terminada \\
  HU-03 & Como equipo, quiero preparar la presentación del taller, para explicar cada pantalla en orden. & Media & 2 & Terminada \\
  \bottomrule
\end{tabularx}

% ---------- 4. Desglose de tareas ----------
\section{Desglose de tareas}

\begin{tabularx}{\textwidth}{c X l c c}
  \toprule
  \textbf{HU} & \textbf{Tarea} & \textbf{Responsable} & \textbf{Horas est.} & \textbf{Estado} \\
  \midrule
  HU-01 & Dibujar a mano las cinco pantallas & Beiker Jaramillo & 5 & Terminada \\
  HU-01 & Definir el flujo y el orden de las pantallas & Laura Ortiz & 2 & Terminada \\
  HU-02 & Escribir el texto que explica cada pantalla & Sarahy Silva & 2 & Terminada \\
  HU-02 & Revisar el prototipo contra la propuesta de solución & Miguel Gómez & 1 & Terminada \\
  HU-03 & Armar la presentación del taller & Valentina Carrillo & 2 & Terminada \\
  \bottomrule
\end{tabularx}

% ---------- 5. Seguimiento diario (Daily Scrum) ----------
\section{Seguimiento diario (Daily Scrum)}

\begin{tabularx}{\textwidth}{c l X X X}
  \toprule
  \textbf{Día} & \textbf{Miembro} & \textbf{Ayer} & \textbf{Hoy} & \textbf{Impedimentos} \\
  \midrule
  1 & Laura Ortiz & Definimos las cinco pantallas & Ordenar el flujo de inicio a acción recomendada & Ninguno \\
  2 & Beiker Jaramillo & Dibujamos inicio y diagnóstico & Dibujar resultado, plan y acción recomendada & Ninguno \\
  \bottomrule
\end{tabularx}

\subsection{Impedimentos y acciones}

\begin{tabularx}{\textwidth}{X l l}
  \toprule
  \textbf{Impedimento} & \textbf{Responsable} & \textbf{Estado} \\
  \midrule
  Ninguno en este sprint & Equipo & No aplica \\
  \bottomrule
\end{tabularx}

% ---------- 6. Métricas del sprint ----------
\section{Métricas del sprint}

\begin{tabularx}{\textwidth}{X c c}
  \toprule
  \textbf{Indicador} & \textbf{Planeado} & \textbf{Real} \\
  \midrule
  Puntos de historia completados & 8 & 8 \\
  Historias completadas          & 3 & 3 \\
  Velocidad del equipo           & 8 & 8 \\
  \bottomrule
\end{tabularx}

% ---------- 7. Revisión del sprint ----------
\section{Revisión del sprint}

\subsection{Incrementos entregados}
\begin{itemize}[leftmargin=1.5em]
  \item {Prototipo a mano de cinco pantallas, desde el inicio hasta la acción recomendada.}
  \item {Presentación del taller práctico con la explicación de cada pantalla.}
\end{itemize}

\subsection{Retroalimentación de los interesados}
{Comentarios del profesor sobre el prototipo entregado por Beiker (nota: 50/50, equivalente a 5,0):}
\begin{itemize}[leftmargin=1.5em]
  \item {Buen trabajo.}
  \item {Le gustó que lo hecho a mano se detallara con gráficos y textos complementarios.}
  \item {Se ven varias de las funcionalidades base de la propuesta de solución, lo que es importante para consolidar un MVP final junto con lo presentado por los demás compañeros.}
  \item {La nota incluye el bono por hacerlo a mano.}
\end{itemize}

\subsection{Historias no completadas}
\begin{tabularx}{\textwidth}{c X X}
  \toprule
  \textbf{ID} & \textbf{Motivo} & \textbf{Acción propuesta} \\
  \midrule
  -- & Todas las historias se completaron. & No aplica \\
  \bottomrule
\end{tabularx}

% ---------- 8. Retrospectiva ----------
\section{Retrospectiva del sprint}

\subsection{Qué salió bien}
\begin{itemize}[leftmargin=1.5em]
  \item {El flujo completo se entiende en cinco pantallas.}
  \item {Cada pantalla tiene una explicación corta de su propósito.}
\end{itemize}

\subsection{Qué se puede mejorar}
\begin{itemize}[leftmargin=1.5em]
  \item {El prototipo no se probó todavía con usuarios reales.}
  \item {Falta unir los prototipos de todos en una sola versión para el MVP.}
\end{itemize}

\subsection{Acciones de mejora}
\begin{tabularx}{\textwidth}{X l l}
  \toprule
  \textbf{Acción} & \textbf{Responsable} & \textbf{Fecha límite} \\
  \midrule
  Consolidar los prototipos del equipo en un MVP & Miguel Gómez & 16/11/2026 \\
  Probar el prototipo con cinco jóvenes & Laura Ortiz & 24/11/2026 \\
  \bottomrule
\end{tabularx}

% ---------- 9. Aprobaciones ----------
\section{Aprobaciones}

\vspace{1cm}
\begin{tabularx}{\textwidth}{X X}
  \textbf{Product Owner} & \textbf{Scrum Master} \\[2cm]
  \rule{6cm}{0.4pt} & \rule{6cm}{0.4pt} \\
  Firma y fecha & Firma y fecha \\
\end{tabularx}

\end{document}
